\chapter{Emergency Medicine}

\medterm{Abscess}-- A localized collection of pus within a tissue, organ, or cavity, typically caused by bacterial infection. Abscesses form when bacteria (most commonly Staphylococcus aureus including MRSA) invade tissue, triggering an intense inflammatory response with neutrophil accumulation and tissue necrosis. The body attempts to wall off the infection, creating a walled-off cavity filled with purulent material (dead neutrophils, bacteria, necrotic tissue). Clinical presentation includes a painful, swollen, erythematous, warm mass that may fluctuate (indicating fluid collection). Systemic symptoms (fever, leukocytosis) may be present. Diagnosis is primarily clinical; ultrasound confirms the presence and size of a fluid collection. Treatment requires incision and drainage (I and D) for Source control, as antibiotics alone cannot penetrate the abscess wall effectively. Small abscesses (< 2 cm) may respond to warm compresses and antibiotics. Culture should be obtained to guide antibiotic therapy, particularly for recurrent abscesses or MRSA risk factors. Complications include cellulitis, bacteremia, and sepsis if untreated. Common locations include skin and soft tissue (furuncles, carbuncles), perianal (perianal abscess), intra-abdominal (from perforated viscus), dental (dental abscess), and brain (cerebral abscess).

\medterm{Acute Angle-Closure Glaucoma}-- Ophthalmologic emergency. Elevated intraocular pressure causing optic nerve damage. Symptoms: severe eye pain, headache, nausea, vomiting, halos around lights, blurred vision. Signs: fixed mid-dilated pupil, corneal edema, rock-hard eye. Treatment: medical (acetazolamide, pilocarpine, mannitol), laser peripheral iridotomy.

\medterm{Acute Aortic Syndromes}-- A spectrum of life-threatening conditions involving the thoracic aorta, including aortic dissection, intramural hematoma (IMH), and penetrating aortic ulcer (PAU). Aortic dissection involves an intimal tear creating a false lumen. Classification (Stanford): Type A involves the ascending aorta (surgical emergency); Type B involves only the descending aorta (medical management unless complicated). Clinical presentation: sudden onset of severe, tearing or ripping chest or back pain radiating to the interscapular region, hypertension or hypotension, pulse deficits, aortic regurgitation murmur, or neurological deficits. Diagnosis: CT angiography (gold standard), transesophageal echocardiography (TEE), or MRI. Management: anti-impulse therapy (IV beta-blockers like esmolol/labetalol targeting HR < 60 bpm and SBP 100--120 mmHg), followed by emergency surgery for Type A or TEVAR/medical management for Type B.

\medterm{Acute Coronary Syndrome (ACS)}-- A spectrum of conditions caused by reduced coronary blood flow and myocardial ischemia: unstable angina, non-ST-elevation MI (NSTEMI), and ST-elevation MI (STEMI). Usual cause: rupture/erosion of an atherosclerotic plaque with superimposed thrombosis. Features: chest pain (pressure/crushing, radiating to arm/jaw), dyspnea, diaphoresis, nausea; atypical in women, diabetics, elderly. Diagnosis: ECG (ST elevation for STEMI), serial troponin. Management: aspirin, antiplatelet (P2Y12), anticoagulant, nitroglycerin, morphine (as needed), statin, beta-blocker; reperfusion—primary PCI within 90 minutes (STEMI) or fibrinolysis.

\medterm{Acute Coronary Syndromes}-- A spectrum of conditions comprising unstable angina (UA), non-ST-elevation myocardial infarction (NSTEMI), and ST-elevation myocardial infarction (STEMI). Pathophysiology: atherosclerotic plaque rupture or erosion leading to thrombus formation and acute coronary artery occlusion. STEMI: complete occlusion presenting with ST-segment elevation in contiguous leads (>= 1 mm in limb leads, >= 2 mm in precordial leads) or new LBBB. NSTEMI: partial occlusion causing cardiac troponin elevation without ST elevation. Unstable angina: ischemic symptoms at rest or increasing frequency without troponin elevation. Management: STEMI requires immediate reperfusion via primary percutaneous coronary intervention (PCI) within 90 minutes (door-to-balloon) or fibrinolysis within 30 minutes if PCI is unavailable within 120 minutes. Medical therapy: dual antiplatelet therapy (aspirin + P2Y12 inhibitor), anticoagulation (heparin), high-intensity statin, beta-blocker, and ACE inhibitor.

\medterm{Acute Pancreatitis}-- Inflammation of pancreas with potential systemic effects. Diagnosis requires 2 of 3: abdominal pain consistent with disease, serum lipase/amylase >= 3x upper limit of normal, characteristic findings on imaging. Etiologies: gallstones (40%), alcohol (30%), hypertriglyceridemia, ERCP, medications, trauma, idiopathic. Severity scoring: APACHE II, Ranson criteria, BISAP. Management: aggressive fluid resuscitation, pain control, early enteral nutrition, treat underlying cause. Complications: pancreatic necrosis, pseudocyst, organ failure.

\medterm{Acute Respiratory Distress Syndrome (ARDS)}-- Progressive respiratory failure due to diffuse alveolar damage. Diagnostic criteria (Berlin definition): acute onset within 1 week, bilateral opacities on chest imaging, respiratory failure not fully explained by cardiac failure, PaO2/FiO2 ratio <= 300 mmHg. Management: lung-protective ventilation (tidal volume 6 mL/kg predicted body weight), prone positioning, conservative fluid strategy. Mortality: 30-40%.

\medterm{Adrenal Crisis}-- An acute, life-threatening manifestation of severe adrenal insufficiency. Causes: primary (Addison disease), secondary (pituitary failure), or tertiary (exogenous corticosteroid withdrawal). Precipitants: systemic infection, trauma, surgery, or abrupt steroid cessation. Clinical presentation: severe hypotension refractory to vasopressors and fluid resuscitation, fever, abdominal pain, vomiting, hyponatremia, hyperkalemia, hypoglycemia, and altered mental status. Emergency management: immediate administration of IV hydrocortisone 100 mg bolus, followed by 200 mg/24 hours continuous infusion; aggressive isotonic saline fluid resuscitation; dextrose for hypoglycemia; and identification and treatment of the underlying precipitating cause.

\medterm{Advanced Cardiovascular Life Support (ACLS)}-- A set of clinical interventions for the urgent treatment of cardiac arrest, stroke, myocardial infarction, and life-threatening arrhythmias. Includes airway management, ECG rhythm interpretation, electrical defibrillation/cardioversion, transcutaneous pacing, and administration of emergency medications (epinephrine, amiodarone, atropine, adenosine, lidocaine). ACLS algorithms guide management of cardiac arrest, bradyarrhythmias, tachyarrhythmias, and acute coronary syndromes.

\medterm{Airway Management}-- Securing and maintaining a patent airway is the first priority in emergency care. Techniques: jaw thrust, head tilt-chin lift, insertion of oropharyngeal or nasopharyngeal airway, endotracheal intubation, surgical airway (cricothyrotomy). Indications: coma (GCS < 8), airway obstruction, respiratory failure, protected airway during resuscitation. Complications: dental injury, esophageal intubation, aspiration, laryngospasm.

\medterm{Altitude Illness}-- Acute mountain sickness (AMS): headache + 1 symptom (nausea, fatigue, dizziness, insomnia) at >2500m. HACE: AMS + ataxia/altered mental status (cerebral edema). HAPE: dyspnea at rest, cough, frothy sputum, crackles, hypoxemia (noncardiogenic pulmonary edema). Prevention: gradual ascent (300-500m/day above 3000m), acetazolamide 125 mg q12h prophylaxis. Treatment: descent (immediate for HACE/HAPE), oxygen, acetazolamide, dexamethasone 8 mg load then 4 mg q6h (HACE), nifedipine 30 mg ER q12h (HAPE), portable hyperbaric chamber (Gamow bag).

\medterm{Anaphylaxis}-- A severe, potentially life-threatening systemic hypersensitivity reaction characterized by rapid onset of multisystem involvement, primarily affecting the skin, respiratory tract, cardiovascular system, and gastrointestinal tract. It is mediated by IgE-dependent mast cell and basophil degranulation (Type I hypersensitivity) or by direct mast cell activation (non-IgE mediated, e.g., radiocontrast, opioids, vancomycin). Common triggers: foods (peanuts, tree nuts, shellfish, fish, eggs, milk, soy, wheat), medications (penicillins, NSAIDs, neuromuscular blocking agents, radiocontrast), insect stings (Hymenoptera: bees, wasps, fire ants), and exercise (exercise-induced anaphylaxis, often fooddependent). Clinical criteria: acute onset (minutes to hours) with involvement of skin/mucosa PLUS respiratory compromise (wheeze, stridor, dyspnea, hypoxia) OR reduced blood pressure (systolic less than 90 mmHg in adults, or greater than 30 percent decrease from baseline) OR gastrointestinal symptoms (vomiting, crampy abdominal pain) with skin involvement. Management: immediate intramuscular epinephrine (0.3-0.5 mg of 1:1,000 solution in the mid-outer thigh, repeat every 5-15 minutes as needed) is the first-line and most critical intervention. Adjunctive treatments: high-flow oxygen, IV fluids (normal saline or lactated Ringer's for hypotension), H1 antihistamine (diphenhydramine 25-50 mg IV), H2 antihistamine (famotidine 20 mg IV), corticosteroids (methylprednisolone 125 mg IV or prednisone 40-60 mg PO), and albuterol nebulizer for bronchospasm. Patients should be observed for at least 4-6 hours (biphasic reaction occurs in 1-20 percent of cases). All patients discharged after anaphylaxis should receive an epinephrine auto-injector and referral to an allergist.

\medterm{Anaphylaxis}-- See anaphylaxis above for full entry.

\medterm{Aortic Dissection}-- A spectrum of life-threatening conditions involving the thoracic aorta, including aortic dissection, intramural hematoma (IMH), and penetrating aortic ulcer (PAU). See acute aortic syndromes.

\medterm{Appendicitis}-- Inflammation of appendix. Classic presentation: periumbilical pain migrating to right lower quadrant, anorexia, nausea, vomiting, low-grade fever. Diagnosis: clinical (Alvarado score), CT abdomen (sensitivity 94%, specificity 95%). Treatment: appendectomy (laparoscopic preferred), antibiotics for uncomplicated cases. Complications: perforation, abscess, peritonitis.

\medterm{Arrhythmias}-- Abnormal heart rhythms. Tachyarrhythmias: atrial fibrillation, atrial flutter, ventricular tachycardia, supraventricular tachycardia. Bradyarrhythmias: sinus bradycardia, heart blocks. Management: ACLS algorithms, rate/rhythm control, anticoagulation when indicated.

\medterm{Asphyxia}-- A life-threatening lack of oxygen due to drowning, choking, or an obstruction of the airways. Management: remove obstruction, rescue breathing, CPR if needed. Foreign body airway obstruction: Heimlich maneuver (abdominal thrusts) for conscious adults; back blows and chest thrusts for infants. Choking: coughing effective for partial obstruction; complete obstruction requires intervention.

\medterm{Aspiration}-- The inhalation of foreign material (oropharyngeal secretions, gastric contents, food, liquid, or blood) into the larynx and lower respiratory tract. Aspiration pneumonitis (Mendelson syndrome) is a chemical lung injury caused by acidic gastric fluid. Aspiration pneumonia is an infectious lung process developing after inhalation of colonized oral secretions. Risk factors: impaired consciousness (alcohol intoxication, stroke, seizure, general anesthesia), dysphagia, enteral tube feeding, and vomiting. Clinical features: sudden cough, dyspnea, wheezing, hypoxemia, fever, and infiltrate on chest radiograph (commonly right lower lobe due to steeper right mainstem bronchus angle). Management: airway suctioning, supplemental oxygen, bronchodilators, and empiric antibiotics for aspiration pneumonia.

\medterm{Aspiration Pneumonia}-- Infectious pneumonia from aspiration of oropharyngeal secretions. Risk factors: dysphagia, impaired consciousness, enteral feeding. Common pathogens: anaerobes, gram-negative bacilli. Treatment: antibiotics covering oral flora.

\medterm{Aspiration Pneumonitis}-- Chemical injury from gastric acid aspiration. Mendelson syndrome. Rapid onset of hypoxemia, bronchospasm, pulmonary edema. Management: oxygen, bronchodilators, suctioning. Prophylactic antibiotics not recommended.

\medterm{Avulsion}-- An injury in which tissue is forcibly torn from the body by shearing or pulling. Common forms include dental, fingertip, nail, and degloving avulsions. Management focuses on hemorrhage control, debridement, tissue preservation, and urgent consultation. Avulsed teeth should be reimplanted within 60 minutes, handled by the crown, and stored in milk, saline, or saliva.

\medterm{Bag-Mask Ventilation}-- Manual ventilation of a non-breathing or inadequately breathing patient using a self-inflating bag (AMBU) attached to a mask sealed over the face. Effective technique: correct mask size, jaw thrust/head extension, two-person technique (two-hand seal), adequate tidal volumes with ETCO2 monitoring. Indicated in respiratory failure, arrest, or pre-intubation; complications: gastric insufflation, aspiration, barotrauma. Used for brief support and as a bridge to a definitive airway.

\medterm{Basic Life Support (BLS)}-- Level of medical care used for victims of life-threatening illnesses or injuries until full medical care can be provided. Focuses on the High-Quality CPR compression-to-ventilation ratio (30:2 in adults), chest compression depth (2-2.4 inches, 100-120/min), early use of an Automated External Defibrillator (AED), and relief of foreign body airway obstruction.

\medterm{Bell's Palsy}-- Idiopathic facial nerve paralysis. Acute onset unilateral facial weakness. May follow viral infection. Treatment: corticosteroids within 72 hours, eye protection. Most recover spontaneously within 3-6 months.

\medterm{Botulism}-- A rare, potentially fatal neuroparalytic disease caused by botulinum toxin, a potent presynaptic neurotoxin produced by Clostridium botulinum, an anaerobic, spore-forming, gram-positive bacillus. See infectious disease.

\medterm{Burns}-- Tissue damage from thermal, chemical, electrical, or radiation sources. Classification: first-degree (epidermis, erythema), second-degree (partial thickness, blisters), third-degree (full thickness, eschar). Assessment: percentage of body surface area (BSA) using Lund-Browder or "rule of nines." Fluid resuscitation: Parkland formula (4 mL x kg x %BSA, half in first 8 hours). Complications: hypovolemic shock, infection, respiratory compromise (inhalation injury), compartment syndrome. Transfer to burn center for: partial-thickness burns >10% BSA, full-thickness burns, burns involving face/hands/perineum, electrical/chemical burns, inhalation injury, burns in patients with comorbidities.

\medterm{Burns}-- See burns above for full entry.

\medterm{Carbon Monoxide Poisoning}-- Toxicity from inhalation of carbon monoxide (CO), a colorless, odorless gas binding hemoglobin ~200-300 times more avidly than oxygen, forming carboxyhemoglobin (COHb) and reducing oxygen delivery. See toxicology.

\medterm{Cardiac Arrest}-- Cessation of effective cardiac output. Causes: ventricular fibrillation (VF), pulseless ventricular tachycardia (pVT), asystole, pulseless electrical activity (PEA). Management: high-quality CPR, defibrillation for VF/pVT (200 J biphasic), epinephrine 1 mg IV every 3-5 minutes, amiodarone for refractory VF/pVT. Return of spontaneous circulation (ROSC) followed by post-cardiac arrest care: hemodynamic optimization, temperature management, coronary angiography if indicated.

\medterm{Cardiopulmonary Resuscitation (CPR)}-- Emergency procedure combining chest compressions and rescue breaths to maintain circulation and oxygenation during cardiac arrest. Compression depth: 2-2.4 inches (5-6 cm) in adults. Rate: 100-120/min. Allow full chest recoil. Minimize interruptions. Compression-to-ventilation ratio: 30:2 for single rescuer, 15:2 for two rescuers with advanced airway. Early defibrillation improves survival.

\medterm{Cat Scratch Disease}-- A regional lymphadenitis caused by Bartonella henselae, a gramnegative bacillus transmitted to humans through cat scratches, bites, or licks on broken skin, particularly from kittens (which carry higher bacterial loads). See infectious disease.

\medterm{Catalepsy}-- See psychiatry.

\medterm{Cauda Equina Syndrome}-- Compression of lumbar and sacral nerve roots below spinal cord termination. Medical emergency. Symptoms: severe low back pain, saddle anesthesia, bowel/bladder dysfunction (urinary retention, incontinence), lower extremity weakness. Etiologies: herniated disc, tumor, abscess, trauma. Diagnosis: MRI spine. Treatment: emergency surgical decompression within 24-48 hours.

\medterm{Central Line-Associated Bloodstream Infection (CLABSI)}-- Bloodstream infection associated with central venous catheter. Prevention: maximal sterile barriers, chlorhexidine skin preparation, optimal site selection (subclavian preferred), daily review of line necessity. Diagnosis: paired blood cultures (peripheral and line). Treatment: antibiotics, catheter removal if possible.

\medterm{Cerebritis}-- Inflammatory process in brain parenchyma without abscess formation. May precede abscess. Causes: bacterial, fungal, parasitic infections. Diagnosis: MRI with contrast showing enhancement. Treatment: antibiotics/antifungals, surgical drainage if abscess forms.

\medterm{Cerebrovascular Accident (Stroke)}-- Interruption of blood flow to the brain, causing neurological deficits. Types: ischemic (87%, thrombotic or embolic), hemorrhagic (13%, intracerebral or subarachnoid). Ischemic stroke treatment: thrombolysis (alteplase within 4.5 hours), thrombectomy (large vessel occlusion within 24 hours). Hemorrhagic stroke: blood pressure control, reversal of anticoagulation, neurosurgical intervention. Prevention: antiplatelets, statins, blood pressure control, anticoagulation for atrial fibrillation.

\medterm{Chemical Injury}-- Tissue damage from corrosive substances. Acid burns: coagulation necrosis, eschar formation. Alkali burns: liquefaction necrosis, deeper penetration. Management: copious irrigation, remove contaminated clothing, assess for airway involvement.

\medterm{Cholinergic Crisis}-- Excess acetylcholine from organophosphate or carbamate poisoning. SLUDGE syndrome: salivation, lacrimation, urination, defecation, GI upset, emesis. Miosis, bradycardia, bronchospasm, muscle fasciculations. Treatment: atropine, pralidoxime.

\medterm{Cocaine Toxicity}-- Sympathomimetic toxidrome: hypertension, tachycardia, hyperthermia, agitation, seizures, myocardial infarction, stroke. Management: benzodiazepines first-line, cooling, activated charcoal if recent ingestion. Avoid beta-blockers (unopposed alpha stimulation).

\medterm{Compartment Syndrome}-- Increased pressure within fascial compartment compromising circulation. Causes: trauma, fracture, burn, cast too tight. Symptoms: pain out of proportion, pain with passive stretch, paresthesia, pallor, paralysis, pulselessness (late). Diagnosis: compartment pressure measurement (>30 mmHg diagnostic). Treatment: emergent fasciotomy.

\medterm{Compartment Syndrome}-- Increased pressure within fascial compartment compromising circulation. See compartment syndrome.

\medterm{Computed Tomography (CT) Scan}-- Cross-sectional imaging using X-rays and computer processing. Emergency uses: head trauma (hemorrhage, fracture), chest (pulmonary embolism, aortic dissection), abdomen (appendicitis, bowel obstruction), CT pulmonary angiography (CTPA), CT angiography (CTA). Fast, widely available, high sensitivity for acute conditions. Radiation exposure is a concern.

\medterm{Concussion}-- Mild traumatic brain injury causing transient neurological dysfunction. Symptoms: headache, confusion, memory problems, dizziness, nausea, sensitivity to light/noise. Management: physical and cognitive rest, gradual return to activity, symptom-limited exercise. Most recover within 7-10 days.

\medterm{Coronary Artery Disease}-- Narrowing or occlusion of coronary arteries, usually due to atherosclerosis, reducing blood flow to the myocardium. See cardiology.

\medterm{Coup-Contrecoup Injury}-- Brain injury from acceleration-deceleration trauma. Coup: injury at site of impact. Contrecoup: injury on opposite side. Common in falls, motor vehicle accidents.

\medterm{Craniocerebral Injury}-- Trauma to scalp, skull, or brain. See head injury.

\medterm{Critical Care}-- Specialty managing critically ill patients. See intensive care medicine.

\medterm{Crush Syndrome}-- Systemic manifestations of prolonged muscle crush injury (rhabdomyolysis) following trauma — typically from collapsed buildings, landslides, entrapment under debris, or excessive compression of limbs. Bywaters syndrome. Pathophysiology: prolonged muscle ischemia, then reperfusion on release, releases myoglobin, potassium, creatine kinase, lactic acid, and phosphorus. Triad of injury: (1) hyperkalemia (cardiac arrhythmia, the leading cause of early death), (2) hypovolemic shock from third-spacing, (3) myoglobinuric acute kidney injury (pigment cast nephropathy, distal renal tubule necrosis). Clinical: trapped $> 4$ hours, comparts syndrome of crushed limb, dark "tea-colored" urine, oliguria. Pre-hospital: do NOT release limb abruptly, apply tourniquet before extrication, IV fluids (1.5 L/h, normal saline with bicarbonate to alkalinize urine to pH $> 6.5$) before release. Hospital: aggressive fluid, mannitol if hypervolemic, treat hyperkalemia (calcium gluconate, insulin/dextrose, $\beta_2$ agonists, kayexalate), dialysis for established AKI, fasciotomy for compartment syndrome.

\medterm{Decerebrate Posturing}-- Abnormal extension posturing indicating severe brain injury. Arms extended, legs extended, internal rotation. Poor prognosis.

\medterm{Decorticate Posturing}-- Abnormal flexion posturing indicating brain injury. Arms flexed, legs extended. Better prognosis than decerebrate.

\medterm{Delirium}-- Acute onset fluctuating disturbance in attention and awareness. Usually reversible. Precipitants: infection, medications, metabolic derangement, pain, constipation, urinary retention. Differentiate from dementia (chronic, progressive). Treatment: identify and treat underlying cause, environmental modification, avoid anticholinergics.

\medterm{Dengue Fever}-- An acute mosquito-borne viral illness caused by dengue virus (DENV-1, -2, -3, -4), a flavivirus transmitted primarily by Aedes aegypti mosquitoes. See infectious disease.

\medterm{Diabetic Hyperosmolar Hyperglycemic State (HHS)}-- Severe hyperglycemia (>600 mg/dL), hyperosmolality (>320 mOsm/kg), dehydration without significant ketoacidosis. Type 2 diabetes. Mortality: 10-20%. Treatment: aggressive fluid resuscitation, insulin, electrolyte replacement.

\medterm{Diabetic Ketoacidosis (DKA)}-- Life-threatening complication of diabetes mellitus (usually type 1) characterized by hyperglycemia (>250 mg/dL), ketonemia, and metabolic acidosis (pH < 7.3, bicarbonate < 18 mEq/L). Precipitated by infection, missed insulin, new-onset diabetes. Symptoms: polyuria, polydipsia, nausea, vomiting, abdominal pain, Kussmaul respirations, fruity breath. Management: IV fluids (normal saline), insulin infusion, potassium replacement, treat underlying cause. Monitor glucose, ketones, electrolytes, anion gap.

\medterm{Diffuse Axonal Injury}-- Shearing of axons from acceleration-deceleration trauma. Causes immediate coma. Poor prognosis. Diagnosis: MRI shows punctate hemorrhages at gray-white junction, brainstem.

\medterm{Drowning}-- Respiratory impairment from submersion in liquid. Primary drowning: immediate respiratory failure. Secondary drowning: delayed pulmonary edema. Management: rescue breathing, CPR, oxygen, warm fluids, treat hypothermia. Prevention: supervision, barriers, life jackets.

\medterm{Drug Overdose}-- Excessive ingestion of medication. Common: opioids, benzodiazepines, acetaminophen, salicylates, tricyclic antidepressants. Management: ABCs, activated charcoal, specific antidotes (naloxone, flumazenil, N-acetylcysteine), supportive care.

\medterm{Eclampsia}-- See obstetrics\_gynecology.

\medterm{Electrocution}-- Injury from electrical current. Types: low voltage (<1000V), high voltage (>1000V). Effects: cardiac arrhythmias, muscle contractions, burns, neurological damage. Management: ABCs, cardiac monitoring, tetanus prophylaxis, wound care, check for compartment syndrome.

\medterm{Emergency Department}-- Hospital department providing immediate medical care for acute illnesses and injuries. Staffed by emergency physicians and supported by specialists. Key principles: triage (prioritize by severity), stabilization, diagnosis, treatment, disposition. Departments should have protocols for common emergencies: chest pain, stroke, trauma, sepsis, respiratory failure.

\medterm{Encephalitis}-- Inflammation of the brain parenchyma, most commonly caused by viral infections (herpes simplex virus, arboviruses, enteroviruses). See infectious disease.

\medterm{Epiglottitis}-- Inflammation of epiglottis. Medical emergency. Symptoms: sore throat, fever, drooling, dysphagia, stridor, tripod positioning. Diagnosis: lateral neck X-ray (thumb sign), direct visualization. Treatment: secure airway, antibiotics (ceftriaxone, vancomycin).

\medterm{Ethylene Glycol Poisoning}-- Antifreeze poisoning. See toxicology.

\medterm{Eye Trauma}-- Ocular injuries ranging from corneal abrasions to globe rupture. Chemical burns: irrigate immediately with saline or water for 30 minutes. Penetrating injuries: shield eye, no pressure, emergent ophthalmology consultation. Foreign bodies: remove superficial items; embedded objects require surgical removal. Retinal detachment: flashes, floaters, curtain over vision--emergency.

\medterm{Factitious Disorder}-- Presentation of physical or psychological symptoms or imposing injury or disease on another, with the deceptive intention of assuming the sick role. See psychiatry.

\medterm{Fever}-- Elevated body temperature (>38 degrees C or 100.4 degrees F). Workup depends on duration, height, and patient factors. Common causes: infections (bacterial, viral), autoimmune disorders, malignancies, drug reactions. Management: antipyretics (acetaminophen, ibuprofen), treat underlying cause. Febrile seizures: generally benign, prevent recurrence with fever control.

\medterm{Flail Chest}-- Multiple rib fractures creating free-floating chest segment. Paradoxical chest movement. Underlying pulmonary contusion common. Management: pain control, oxygen, ventilation if needed, surgical fixation for severe cases.

\medterm{Foreign Body}-- See foreign body above.

\medterm{Foreign Body Aspiration}-- Inhalation of object into airway. Common in children (peanuts, coins, small toys). Signs: choking, coughing, stridor, wheezing, decreased breath sounds. Diagnosis: chest X-ray (may be negative), bronchoscopy. Management: back blows and chest thrusts for infants, Heimlich maneuver for conscious adults, emergency bronchoscopy for complete obstruction.

\medterm{Frostbite}-- See environmental emergencies.

\medterm{Glasgow Coma Scale (GCS)}-- Standardized scale assessing consciousness. Components: eye opening (1-4), verbal response (1-5), motor response (1-6). Total: 3-15. Severe: 3-8, Moderate: 9-12, Mild: 13-15. Used for head injury assessment.

\medterm{Head Injury}-- Trauma to scalp, skull, or brain. Scalp lacerations: bleed profusely, require suturing. Skull fracture: basilar (Battle sign, raccoon eyes) or depressed (surgical emergency). Concussion: transient neurological dysfunction, manage with rest, gradual return to activity. Intracranial hemorrhage: epidural (lucid interval, arterial bleed), subdural (venous bleed, elderly/alcoholic), subarachnoid (worst headache). CT head is diagnostic.

\medterm{Heat Exhaustion}-- See environmental emergencies.

\medterm{Heat Stroke}-- See environmental emergencies.

\medterm{Hemorrhagic Shock}-- Loss of >20% blood volume (>1L in adult). Signs: tachycardia, hypotension, cool clammy skin, altered mental status, oliguria. Classification: Class I (<15%), Class II (15-30%), Class III (30-40%), Class IV (>40%). Management: control bleeding, IV fluids, blood transfusion (1:1:1 ratio for massive transfusion), vasopressors if needed.

\medterm{Hemothorax}-- Blood in pleural space. Traumatic or spontaneous. Symptoms: chest pain, dyspnea, decreased breath sounds, dullness to percussion. Diagnosis: chest X-ray, ultrasound, CT. Treatment: tube thoracostomy, blood transfusion, surgery if ongoing bleeding.

\medterm{Hepatic Encephalopathy}-- Neuropsychiatric syndrome from liver failure. Symptoms: confusion, asterixis, personality changes, coma. Treatment: lactulose, rifaximin, treat precipitating factors.

\medterm{Hospital-Acquired Pneumonia}-- Pneumonia developing > 48 hours after hospital admission. See infectious disease.

\medterm{Hydrofluoric Acid Exposure}-- See chemical injuries.

\medterm{Hyperbaric Oxygen Therapy}-- Administration of 100% oxygen at increased atmospheric pressure. Indications: carbon monoxide poisoning, decompression sickness, gas gangrene, necrotizing soft tissue infections, delayed radiation injury, crush injury, compromised grafts. Mechanism: increases dissolved oxygen, vasoconstriction, antimicrobial effect.

\medterm{Hyperkalemia}-- Serum potassium >5.0 mEq/L. ECG changes: peaked T waves, widened QRS, sine wave. Causes: renal failure, medications (ACE inhibitors, potassium-sparing diuretics), tissue breakdown, acidosis. Treatment: stabilize myocardium (calcium gluconate), shift potassium (insulin + glucose, albuterol, sodium bicarbonate), remove potassium (loop diuretics, potassium binders, dialysis).

\medterm{Hypoglycemia}-- Blood glucose <70 mg/dL. Symptoms: tremor, palpitations, sweating, hunger, confusion, seizures, coma. Causes: insulin or sulfonylurea overdose, missed meals, alcohol, critical illness. Treatment: conscious patient--oral glucose; unconscious--IV dextrose (D50W) or IM glucagon. Recurrent hypoglycemia requires medication adjustment.

\medterm{Hyponatremia}-- Serum sodium <135 mEq/L. Symptoms: nausea, headache, confusion, seizures, coma. Severity depends on acuity and magnitude. Causes: SIADH, heart failure, cirrhosis, nephrotic syndrome, diuretics, adrenal insufficiency. Treatment: fluid restriction, hypertonic saline (severe/symptomatic), treat underlying cause.

\medterm{Hypothermia}-- Core temperature <35 degrees C (95 degrees F). Stages: mild (32-35), moderate (28-32), severe (<28). J-wave (Osborn wave) on ECG. Management: passive rewarming (mild), active external rewarming (moderate), active internal rewarming (severe). Avoid rough handling (risk of VF).

\medterm{Idiopathic Intracranial Hypertension}-- Elevated intracranial pressure without mass lesion. Symptoms: headache, pulsatile tinnitus, visual changes, papilledema. Treatment: acetazolamide, weight loss, surgical options.

\medterm{Inborn Error of Metabolism}-- See biochemistry.

\medterm{Infection}-- Measures to prevent healthcare-associated infections: hand hygiene, personal protective equipment (PPE), isolation precautions (contact, droplet, airborne), sterilization, disinfection. See infection control.

\medterm{Infection Control}-- Measures to prevent healthcare-associated infections: hand hygiene, personal protective equipment (PPE), isolation precautions (contact, droplet, airborne), sterilization, disinfection. Multidrug-resistant organisms: MRSA, VRE, C. difficile, carbapenem-resistant Enterobacteriaceae.

\medterm{Intensive Care Medicine}-- Specialty managing critically ill patients. Components: monitoring, ventilatory support, hemodynamic support, infection control, nutrition, sedation, pain management. Common conditions: sepsis, ARDS, traumatic brain injury, post-cardiac arrest, multiorgan failure.

\medterm{Juvenile Rheumatoid Arthritis}-- See rheumatology.

\medterm{Kussmaul Respirations}-- Deep, labored breathing pattern associated with diabetic ketoacidosis. Compensatory hyperventilation to blow off CO2 and correct metabolic acidosis. Name: Adolf Kussmaul, German physician.

\medterm{Lactic Acidosis}-- Lactate >4 mmol/L with acidemia. Types: Type A (tissue hypoxia: shock, sepsis, hypoxia), Type B (non-hypoxic: metformin, liver failure, malignancy, drugs). Treatment: address underlying cause, supportive care.

\medterm{Leptospirosis}-- A zoonotic bacterial infection caused by spirochaetes of the genus Leptospira, transmitted through contact with water or soil contaminated with the urine of infected animals (rodents, cattle, dogs). See infectious disease.

\medterm{Lyme Disease}-- A tick-borne infection caused by Borrelia burgdorferi (Ixodes scapularis, deer tick), endemic to the northeastern and upper midwestern US. See infectious disease.

\medterm{Malingering}-- The intentional fabrication or exaggeration of physical or psychological symptoms for external incentives (avoiding work, obtaining financial compensation, evading criminal prosecution, obtaining drugs). See psychiatry.

\medterm{Mass Casualty Incident}-- Event overwhelming available resources, requiring triage and resource allocation. Triage systems: START (Simple Triage and Rapid Treatment), SALT. Categories: immediate (red), delayed (yellow), minor (green), expectant (black). Incident command system (ICS) coordinates response.

\medterm{Mediastinitis}-- Inflammation of mediastinum. Acute: post-sternotomy infection (serious, mortality 10-50%). Chronic: mediastinal masses, granulomatous diseases. Diagnosis: CT chest, cultures. Treatment: surgical debridement, antibiotics.

\medterm{Meningitis}-- Inflammation of the meninges most commonly caused by bacterial infection (Neisseria meningitidis, Streptococcus pneumoniae, Haemophilus influenzae, Group B Streptococcus). See infectious disease.

\medterm{Methemoglobinemia}-- Elevated methemoglobin (>1%) causing functional anemia. Cyanosis unresponsive to oxygen. Causes: dapsone, benzocaine, nitrates, sulfonamides. Diagnosis: co-oximetry, chocolate-colored blood. Treatment: methylene blue 1-2 mg/kg IV, vitamin C alternative.

\medterm{Migraine with Aura}-- A primary headache disorder characterized by recurrent attacks of fully reversible visual, sensory, or speech symptoms that develop gradually over at least five minutes and last no more than sixty minutes, followed by headache with the characteristics of migraine. See neurology.

\medterm{Myocardial Infarction (MI)}-- Death of myocardium due to prolonged ischemia. Types: STEMI (ST elevation), NSTEMI (troponin elevation without ST elevation). Diagnosis: ECG changes, troponin elevation. Treatment: MONA (morphine, oxygen, nitroglycerin, aspirin), reperfusion (PCI or fibrinolysis). Complications: arrhythmias, heart failure, cardiogenic shock, ventricular septal rupture, papillary muscle rupture.

\medterm{Myxedema Coma}-- Severe hypothyroidism with altered mental status. Precipitated by infection, cold exposure, medications. Symptoms: hypothermia, bradycardia, hypotension, hypoventilation, hyponatremia, hypoglycemia. Treatment: IV levothyroxine, hydrocortisone, supportive care. Mortality: 20-60%.

\medterm{Narrow-Complex Tachycardia}-- See cardiology.

\medterm{Necrotizing Fasciitis}-- Rapidly progressive soft tissue infection with fascial necrosis. "Flesh-eating disease." Types: Type I (polymicrobial), Type II (Group A Streptococcus). Symptoms: severe pain out of proportion, erythema, swelling, bullae, crepitus, systemic toxicity. Diagnosis: clinical, imaging (gas in tissues), surgical exploration. Treatment: emergent surgical debridement, broad-spectrum antibiotics.

\medterm{Obstructive Sleep Apnea}-- Recurrent upper airway collapse during sleep. Symptoms: snoring, witnessed apneas, daytime sleepiness, morning headaches. Risk factors: obesity, male sex, older age, neck circumference >17 inches. Diagnosis: polysomnography. Treatment: CPAP, weight loss, oral appliances, surgery.

\medterm{Opioid Overdose}-- A medical emergency characterized by the classic triad of respiratory depression, miosis (pinpoint pupils), and decreased level of consciousness. See toxicology.

\medterm{Organophosphate Poisoning}-- Excess acetylcholine from organophosphate or carbamate poisoning. See cholinergic crisis.

\medterm{Pain}-- Emergency pain control: NSAIDs, opioids, ketamine, lidocaine. See pain management.

\medterm{Pain Management}-- Emergency pain control: NSAIDs, opioids, ketamine, lidocaine. Assessment: location, intensity, character, timing, aggravating/alleviating factors. Pain is the fifth vital sign. Multimodal approach reduces opioid requirements. Neuropathic pain: gabapentin, pregabalin.

\medterm{Paracetamol Toxicity}-- See toxicology.

\medterm{Pediatric Fever}-- See pediatrics.

\medterm{Pericardial Tamponade}-- Fluid accumulation in pericardial space compromising cardiac filling. Beck triad: hypotension, jugular venous distension, muffled heart sounds. Electrical alternans on ECG. Diagnosis: echocardiography. Treatment: emergent pericardiocentesis, pericardial window.

\medterm{Phobias}-- See psychiatry.

\medterm{Plague}-- A bacterial infection caused by Yersinia pestis, transmitted by fleas from rodents. See infectious disease.

\medterm{Post-Concussion Syndrome}-- Persistent symptoms after concussion: headache, dizziness, cognitive difficulties, sleep disturbance, mood changes. Management: symptom-directed, gradual return to activity.

\medterm{Post-Traumatic Stress Disorder}-- See psychiatry.

\medterm{Preeclampsia}-- See obstetrics\_gynecology.

\medterm{Propofol Infusion Syndrome}-- Rare but fatal complication of prolonged high-dose propofol infusion. Features: metabolic acidosis, rhabdomyolysis, cardiac failure, hyperlipidemia. Risk: critical illness, children, high doses (>4-5 mg/kg/hr), prolonged infusion (>48 hours). Prevention: limit dose and duration.

\medterm{Rabies}-- A fatal viral encephalitis caused by rabies virus, transmitted through the bite of infected animals. See infectious disease.

\medterm{Reperfusion Injury}-- Tissue damage following restoration of blood flow after ischemia. Mechanisms: oxidative stress, calcium overload, inflammation. Examples: myocardial reperfusion injury, stroke reperfusion injury. Mitigation: antioxidants, reperfusion protocols.

\medterm{Respiratory Distress}-- Difficulty breathing due to various causes. Assessment: work of breathing (retractions, nasal flaring, tripod position), oxygen saturation, breath sounds. Causes: asthma, COPD, pneumonia, pulmonary embolism, pneumothorax, anaphylaxis, heart failure. Treatment: oxygen, bronchodilators, corticosteroids, antibiotics, specific interventions (needle decompression for tension pneumothorax).

\medterm{Ruptured Abdominal Aortic Aneurysm}-- See vascular surgery.

\medterm{Salicylate Toxicity}-- See toxicology.

\medterm{Sepsis}-- Life-threatening organ dysfunction caused by dysregulated host response to infection. Sepsis-3 definition: infection + SOFA score increase >= 2. Septic shock: sepsis with vasopressor requirement and lactate >2 mmol/L. Management: early antibiotics (within 1 hour), fluid resuscitation (30 mL/kg crystalloid), source control, vasopressors if needed. Procalcitonin may guide antibiotic duration.

\medterm{Serotonin Syndrome}-- A potentially life-threatening condition caused by excessive serotonergic activity in the central and peripheral nervous systems, typically resulting from drug interactions or overdose involving serotonergic medications. See pharmacology.

\medterm{Shock}-- See hemorrhagic shock, septic shock.

\medterm{Spinal Precautions}-- Immobilization of spine in trauma patients to prevent neurological injury. Indications: mechanism (fall >3 feet, MVC, gunshot), altered mental status, focal neurological deficit, midline spinal tenderness. Technique: cervical collar, backboard, log roll. Clearance: NEXUS criteria or Canadian C-spine rule.

\medterm{Spinal Shock}-- Temporary loss of all neural activity below spinal cord injury. Flaccid paralysis, areflexia, loss of sensation. Resolves over days-weeks. Differentiate from neurogenic shock (hemodynamic instability).

\medterm{Spontaneous Pneumothorax}-- Life-threatening pneumothorax with mediastinal shift. See tension pneumothorax.

\medterm{Status Epilepticus}-- Continuous seizure >5 minutes or recurrent seizures without recovery. Medical emergency. Treatment: benzodiazepines (lorazepam, diazepam), then fosphenytoin, valproate, levetiracetam. Refractory: anesthesia (propofol, midazolam).

\medterm{Stress Ulcer Prophylaxis}-- Prevention of stress-related mucosal disease in critically ill patients. Risk factors: mechanical ventilation >48 hours, coagulopathy, trauma, burns, sepsis. Agents: H2 blockers, proton pump inhibitors. Duration: while risk factors present.

\medterm{Sudden Cardiac Death}-- Unexpected death within 1 hour of symptom onset. Most common cause: ventricular fibrillation. Risk factors: coronary artery disease, cardiomyopathy, channelopathies. Prevention: ICD implantation, antiarrhythmics.

\medterm{Tension Pneumothorax}-- Life-threatening pneumothorax with mediastinal shift. Emergency needle decompression (2nd intercostal space, midclavicular line) followed by chest tube. Signs: respiratory distress, hypotension, tracheal deviation, distended neck veins, unilateral absent breath sounds.

\medterm{Tetanus}-- An acute neurological disease caused by tetanospasmin toxin from Clostridium tetani, introduced through wounds. See infectious disease.

\medterm{Thermal Injury (Burns)}-- Tissue damage from heat. Management: stop burning, remove clothing, cool with water (not ice), cover with sterile dressing. Assess depth and extent. Fluid resuscitation for >15% BSA burns. Tetanus prophylaxis, pain control, wound care.

\medterm{Thyroid Storm}-- Life-threatening exacerbation of hyperthyroidism. Precipitated by infection, surgery, trauma. Symptoms: fever, tachycardia, agitation, delirium, vomiting, diarrhea. Treatment: beta-blockers, thionamides, iodine, corticosteroids, supportive care.

\medterm{Toxic Encephalopathy}-- See toxicology.

\medterm{Toxic Shock Syndrome}-- An acute, life-threatening illness caused by bacterial exotoxins. See infectious disease.

\medterm{Trauma}-- Physical injury from external force. Primary survey: ABCDE (Airway, Breathing, Circulation, Disability, Exposure). Secondary survey: head-to-toe examination, history (AMPLE), diagnostic tests. Injuries: fractures, dislocations, internal bleeding, organ damage. ATLS protocol guides management.

\medterm{Trauma Center}-- Hospital designated to provide comprehensive care for injured patients. Levels: Level I (complete care, research, prevention), Level II (similar but limited resources), Level III (stabilization, transfer), Level IV (advanced trauma life support, transfer). Transfer criteria: need for higher level of care, specific injuries.

\medterm{Treatment-Resistant Depression}-- See psychiatry.

\medterm{Unresponsive Coma}-- See coma.

\medterm{Upper Gastrointestinal Bleeding}-- Bleeding from esophagus, stomach, or duodenum. Causes: peptic ulcer, varices, Mallory-Weiss tear, gastritis. Presentation: hematemesis, melena, hematochezia. Management: ABCs, IV access, type and crossmatch, PPI infusion, endoscopy within 24 hours. Variceal bleeding: octreotide, antibiotics, band ligation.

\medterm{Uremic Pericarditis}-- Pericarditis associated with severe renal failure. Chest pain, pericardial friction rub, diffuse ST elevation. Risk of tamponade. Treatment: intensify dialysis, NSAIDs if not contraindicated.

\medterm{Vasovagal Syncope}-- Most common cause of syncope. Triggered by emotional stress, pain, prolonged standing. Mechanism: vagal stimulation causing bradycardia and vasodilation. Prodrome: nausea, diaphoresis, lightheadedness, tunnel vision. Treatment: reassurance, avoid triggers, physical counterpressure maneuvers.

\medterm{Ventricular Fibrillation (VF)}-- Chaotic electrical activity in ventricles, causing no cardiac output. Immediate defibrillation is treatment. Causes: coronary artery disease, cardiomyopathy, electrolyte disturbances, drug toxicity. Post-resuscitation care: targeted temperature management, coronary angiography, ICU admission.

\medterm{Ventricular Tachycardia}-- A wide-complex tachycardia (QRS >0.12 s) originating from the ventricles, with a rate of 100-250 bpm. See cardiology.

\medterm{Wernicke's Encephalopathy}-- See toxicology.

\medterm{Wound Care}-- Management of traumatic wounds: irrigation, debridement, closure (sutures, staples, glue), dressing. Tetanus prophylaxis based on vaccination history and wound type. Antibiotics for contaminated wounds, animal/human bites, or signs of infection. Healing: primary, secondary, tertiary intention.

\medterm{Wound Infection}-- See infection.

\medterm{Xylitol Toxicity}-- See toxicology.

\medterm{Yellow Fever}-- A mosquito-borne viral hemorrhagic fever caused by yellow fever virus (Flavivirus), endemic in tropical Africa and South America. See infectious disease.

\medterm{Zika Virus}-- A positive-sense RNA virus belonging to the Flaviviridae family that causes mild febrile illness. See infectious disease.

